% 04-nlt.tex — NLT（非线性变换，Non-Linear Transform）
% Source: chapters/04-nlt-nonlinear-transform.md

\chapter{NLT（非线性变换）}\label{ch:nlt}

\section{本章目标}

\begin{itemize}
  \item 理解 NLT 的作用：将无符号像素值映射到以零为中心的有符号系数范围
  \item 掌握三种 NLT 模式的数学定义：None、Quadratic、Extended
  \item 理解 NLT Quadratic 中定点平方根近似算法的工作原理
  \item 理解 NLT Extended 的三段式结构（低区、线性区、高区）
\end{itemize}

\section{背景与标准语义}

NLT（Non-Linear Transform）是 JPEG XS 编码管线的第一个处理步骤。它的作用是将输入的无符号整数像素值 $[0, 2^d - 1]$ 转换为以零为中心的有符号系数，供后续的 MCT 和 DWT 使用。

三种模式的设计意图：

\begin{itemize}
  \item \textbf{None}：线性映射，仅做 DC 电平移除和位宽扩展。适合已经以零为中心的数据或需要 Mathematically Lossless 的场景。
  \item \textbf{Quadratic}：二次映射，$\sqrt{x}$ 前向、$x^2$ 逆向。对暗部区域分配更多精度（压缩动态范围），适合自然图像。
  \item \textbf{Extended}：三段式映射，低区和高区使用二次曲线，中间线性区保持斜率。比 Quadratic 更灵活，可独立控制阈值和斜率。
\end{itemize}

\section{数学定义}

\subsection{通用参数}

\begin{table}[H]
\centering
\caption{NLT 通用参数}
\label{tab:nlt-params}
\begin{tabular}{@{}lp{4.5cm}ll@{}}
\toprule
符号 & 含义（在 NLT 中的作用） & 代码 & 典型值 \\
\midrule
$d$ & 输入图像的像素位深。例如 8-bit 图像 $d=8$，像素范围 $[0,255]$ & \texttt{im->depth} & 8--16 \\
$B_w$ & NLT 输出后的内部位宽。通常 $B_w > d$，多出的位用于后续 DWT 和量化的精度 & \texttt{p.Bw} & 18--20 \\
$s = B_w - d$ & 左移位数。NLT 第一步先把像素值左移 $s$ 位，从 $d$ bit 扩展到 $B_w$ bit & 计算值 & 0--12 \\
$\mathit{dclev} = 2^{B_w-1}$ & DC 电平。左移后的值减去 $\mathit{dclev}$，把无符号范围 $[0,\ldots]$ 平移到以零为中心 & 计算值 & --- \\
$\mathit{max\_val} = 2^d - 1$ & 原始像素的最大可能值。用于 clip 操作的上界 & 计算值 & --- \\
\bottomrule
\end{tabular}
\end{table}

\subsection{NLT None（线性）}

前向变换（编码端）：
\begin{equation}\label{eq:nlt-none-fwd}
v_{\mathit{out}} = (x \ll s) - \mathit{dclev}
\end{equation}

这个式子将 $[0, 2^d-1]$ 映射到 $[-2^{B_w-1}, 2^{B_w-1} - 2^s]$，即以零为中心。

逆向变换（解码端）：
\begin{equation}\label{eq:nlt-none-inv}
x_{\mathit{rec}} = \mathrm{clip}\!\left(\frac{v_{\mathit{out}} + \mathit{dclev} + 2^{s-1}}{2^s},\; 0,\; \mathit{max\_val}\right)
\end{equation}

其中 $2^{s-1}$ 是舍入项。

\subsection{NLT Quadratic（二次）}

\subsubsection{辅助参数}

\begin{table}[H]
\centering
\caption{NLT Quadratic 辅助参数}
\label{tab:nlt-quad-params}
\begin{tabular}{@{}lp{3.5cm}l@{}}
\toprule
符号 & 定义（控制什么） & 代码 \\
\midrule
$\sigma$ & 符号标志（1 bit）。决定 $v_{\mathit{dco}}$ 的偏移方向：0 = 正偏移，1 = 负偏移 & \texttt{quadratic.sigma} \\
$\alpha$ & DC 偏移量（15 bits）。控制零点附近的映射行为，调整暗部压缩曲线 & \texttt{quadratic.alpha} \\
$v_{\mathit{dco}}$ & 等效 DC 电平。NLT 先减去此值再平方根，使映射曲线从合适的起点开始 & $\alpha - \sigma \times 32768$ \\
\bottomrule
\end{tabular}
\end{table}

\subsubsection{前向变换}

\begin{align}
v_1 &= \mathrm{clip}(x - v_{\mathit{dco}},\; 0,\; \mathit{max\_val}) \label{eq:nlt-q-v1} \\
v_2 &= v_1 \ll s \label{eq:nlt-q-v2} \\
v_3 &= \mathrm{sqrt\_approx}(v_2,\; B_w) \label{eq:nlt-q-v3} \\
v_{\mathit{out}} &= v_3 - \mathit{dclev} \label{eq:nlt-q-out}
\end{align}

其中 $\mathrm{sqrt\_approx}(v, B_w)$ 是定点平方根近似算法（见下节）。

这个变换的效果：对大值做平方根压缩，使暗部（小值）保留更多码空间。

\subsubsection{逆向变换}

\begin{align}
v &= \mathrm{clip}(v_{\mathit{out}} + \mathit{dclev},\; 0,\; 2^{B_w} - 1) \label{eq:nlt-q-inv-v} \\
v_{\mathit{rec}} &= v^2 \gg (2B_w - d) \quad \text{with rounding} \label{eq:nlt-q-inv-vrec} \\
x_{\mathit{rec}} &= \mathrm{clip}(v_{\mathit{rec}} + v_{\mathit{dco}},\; 0,\; \mathit{max\_val}) \label{eq:nlt-q-inv-x}
\end{align}

注意前向用 $\sqrt{\cdot}$，逆向用 $(\cdot)^2$。由于定点近似，正逆变换不是严格互逆的。

\subsubsection{定点平方根算法}

\begin{lstlisting}[language=C,caption={定点平方根近似——\texttt{sqrt\_approx\_fixpoint()}（\texttt{nlt.c:85}）}]
static INLINE int64_t sqrt_approx_fixpoint(int64_t v, const uint8_t Bw)
{
    assert(v >= 0);
    int64_t r = 0;
    for (int i = 0; i < Bw; ++i)
    {
        r <<= 1;
        v <<= 2;
        if ((v >> Bw) > r)
        {
            v -= (r + 1) << Bw;
            r += 2;
        }
    }
    return r >> 1;
}
\end{lstlisting}

这是一个逐位（digit-by-digit）的定点平方根算法。每轮迭代处理 1 bit 输出精度，共 $B_w$ 轮，输出精度为 $B_w/2$ bit。

算法本质：
\begin{itemize}
  \item $v$ 被左移 2 位（相当于处理输入的 2 bit）
  \item $r$ 是当前部分结果
  \item 比较 $(v \gg B_w)$ 与 $r$ 决定当前位是 0 还是 1
  \item 共 $B_w$ 轮迭代
\end{itemize}

\subsection{NLT Extended（扩展）}

\subsubsection{参数}

\begin{table}[H]
\centering
\caption{NLT Extended 参数}
\label{tab:nlt-ext-params}
\begin{tabular}{@{}lp{5.5cm}l@{}}
\toprule
符号 & 含义（控制什么） & 代码 \\
\midrule
$T_1$ & 低区→线性区的分界值。像素值低于 $T_1$ 时走二次曲线（暗部压缩），高于 $T_1$ 走线性区 & \texttt{extended.T1} \\
$T_2$ & 线性区→高区的分界值。像素值高于 $T_2$ 时走二次曲线（高光压缩），$T_1$ 到 $T_2$ 之间为线性 & \texttt{extended.T2} \\
$E$ & 线性区的斜率控制参数（3 bit，值 0--7）。$E$ 越大，线性区斜率越陡 & \texttt{extended.E} \\
$e = B_w - E$ & 实际参与运算的辅助指数，$E$ 越大则 $e$ 越小，线性区越平缓 & 计算值 \\
\bottomrule
\end{tabular}
\end{table}

\subsubsection{派生常量}

\begin{align}
B_2 &= T_1^2 \\
A_1 &= B_2 + T_1 \cdot 2^e + 2^{2e-2} \\
B_1 &= T_1 + 2^{e-1} \\
A_3 &= B_2 + T_2 \cdot 2^e - 2^{2e-2} \\
B_3 &= T_2 - 2^{e-1} \\
Q_1 &= B_2 + T_1 \cdot 2^e \\
Q_2 &= B_2 + T_2 \cdot 2^e
\end{align}

\subsubsection{前向变换}

\begin{equation}
v = x \ll s
\end{equation}

三段处理：

\textbf{低区}（$v < Q_1$）：
\begin{equation}
v_{\mathit{out}} = B_1 - \sqrt{A_1 - v} - \mathit{dclev}
\end{equation}

\textbf{线性区}（$Q_1 \leq v < Q_2$）：
\begin{equation}
v_{\mathit{out}} = v - B_2 \,/\, 2^e - \mathit{dclev}
\end{equation}

\textbf{高区}（$v \geq Q_2$）：
\begin{equation}
v_{\mathit{out}} = B_3 + \sqrt{v - A_3} - \mathit{dclev}
\end{equation}

注意：前向变换中低区和高区使用 \texttt{sqrt()}（C 标准库浮点函数），线性区使用整数右移。

\subsubsection{逆向变换}

\begin{equation}
v = v_{\mathit{out}} + \mathit{dclev}
\end{equation}

\textbf{低区}（$v < T_1$）：
\begin{equation}
v_{\mathit{rec}} = \mathrm{clip}\!\left(\frac{A_1 - (B_1 - v)^2 + \mathit{rounding}}{2^{2s}},\; 0,\; \mathit{max\_val}\right)
\end{equation}

\textbf{线性区}（$T_1 \leq v < T_2$）：
\begin{equation}
v_{\mathit{rec}} = \frac{v \cdot 2^e + B_2 + \mathit{rounding}}{2^{2s}}
\end{equation}

\textbf{高区}（$v \geq T_2$）：
\begin{equation}
v_{\mathit{rec}} = \mathrm{clip}\!\left(\frac{(v - B_3)^2 + A_3 + \mathit{rounding}}{2^{2s}},\; 0,\; \mathit{max\_val}\right)
\end{equation}

注意：逆变换中低区和高区使用 $(\cdot)^2$（整数平方），与前向的 $\sqrt{\cdot}$ 互为逆操作。

\subsection{阈值自动生成}

\texttt{xs\_config\_nlt\_extended\_auto\_thresholds()}（\texttt{nlt.c:316}）根据目标 bpp 自动计算 $T_1, T_2$：
\begin{align}
T_1 &= \sqrt{2^{2e-2} + \mathit{th1} \cdot 2^s} - 2^{e-1} \\
T_2 &= \mathit{th2} \cdot 2^s - T_1^2 \,/\, 2^e
\end{align}

其中 $\mathit{th1}$ 和 $\mathit{th2}$ 是用户指定的目标像素值阈值。

\section{处理流程}

\subsection{编码端（前向）}

\begin{lstlisting}[language=C,numbers=none]
for each component c:
    for each sample in component c:
        1. 根据 Tnlt 类型选择变换函数
        2. 应用前向变换公式
        3. 结果存回原位置（in-place）
\end{lstlisting}

\subsection{解码端（逆向）}

\begin{lstlisting}[language=C,numbers=none]
for each component c:
    for each sample in component c:
        1. 根据 Tnlt 类型选择逆变换函数
        2. 应用逆变换公式
        3. 结果存回原位置（in-place）
\end{lstlisting}

注意：NLT 对每个分量独立处理，不涉及分量间交互。

\subsection{流程图}

\begin{figure}[H]
\centering
\begin{tikzpicture}[
    node distance=1cm and 1.8cm,
    block/.style={rectangle, draw=structurecolor!60, fill=structurecolor!10,
                  text width=2.8cm, align=center, minimum height=0.85cm,
                  font=\small, inner sep=4pt},
    decision/.style={diamond, draw=second!70, fill=second!12,
                     text width=3.2cm, align=center, aspect=2.5,
                     font=\small, inner sep=2pt},
    arrow/.style={-{Stealth[length=2.5mm]}, thick, draw=structurecolor!60}
]
    \node[block] (input) {输入像素值};
    \node[decision, below=of input] (type) {NLT 类型?\\$T_{\mathit{nlt}}$};
    \node[block, below left=1.8cm and 2.2cm of type] (none)  {NONE:\\直接传递};
    \node[block, below=2cm of type] (quad) {QUADRATIC:\\定点平方根近似};
    \node[block, below right=1.8cm and 2.2cm of type] (ext) {EXTENDED:\\三段映射};
    \node[block, below=4cm of type] (output) {输出到 MCT};

    \draw[arrow] (input) -- (type);
    \draw[arrow] (type) -| node[above left,font=\footnotesize,pos=0.3] {NONE} (none);
    \draw[arrow] (type) -- node[right,font=\footnotesize,pos=0.4] {QUAD} (quad);
    \draw[arrow] (type) -| node[above right,font=\footnotesize,pos=0.3] {EXT} (ext);
    \draw[arrow] (none) |- (output);
    \draw[arrow] (quad) -- (output);
    \draw[arrow] (ext) |- (output);
\end{tikzpicture}
\caption{NLT 模式选择流程图}
\label{fig:nlt-flowchart}
\end{figure}

\section{实现映射}

\begin{implmap}
\begin{tabular}{@{}L{2.8cm}L{1.5cm}L{4.0cm}L{4.5cm}@{}}
\toprule
标准概念 & 入口文件 & 关键函数/结构 & 说明 \\
\midrule
NLT 前向总入口 & \btt{nlt.c} & \btt{nlt\_forward\_transform()} & 根据 \btt{Tnlt} 分发到具体函数 \\
NLT 逆向总入口 & \btt{nlt.c} & \btt{nlt\_inverse\_transform()} & 根据 \btt{Tnlt} 分发 \\
线性前向 & \btt{nlt.c:121} & \btt{nlt\_forward\_linear()} & Eq~\eqref{eq:nlt-none-fwd} \\
二次前向 & \btt{nlt.c:162} & \btt{nlt\_forward\_quadratic()} & Eq~\eqref{eq:nlt-q-v1}--\eqref{eq:nlt-q-out} \\
扩展前向 & \btt{nlt.c:226} & \btt{nlt\_forward\_extended()} & 三段式变换 \\
定点平方根 & \btt{nlt.c:85} & \btt{sqrt\_approx\_fixpoint()} & 逐位逼近算法 \\
阈值自动生成 & \btt{nlt.c:316} & \btt{xs\_config\_nlt\_extended\_auto\_thresholds()} & 自动计算 $T_1, T_2$ \\
\bottomrule
\end{tabular}
\end{implmap}

\begin{codefact}
本章关键结论依据以下函数：
\begin{itemize}
  \item \texttt{nlt\_forward\_transform()} --- \texttt{nlt.c} --- NLT 正变换入口，根据 Tnlt 分发
  \item \texttt{nlt\_forward\_quadratic()} --- \texttt{nlt.c:162} --- 二次变换核心
  \item \texttt{sqrt\_approx\_fixpoint()} --- \texttt{nlt.c:85} --- 定点平方根近似
\end{itemize}
\end{codefact}

\section{常见误解}

\begin{misunderstanding}
\begin{enumerate}
  \item \textbf{``NLT 是有损的''} --- NLT None 在整数域内是无损的。Quadratic 和 Extended 由于定点近似引入微小误差，但这通常比后续量化引入的误差小得多。
  \item \textbf{``Quadratic 和 Extended 的平方/平方根是精确的''} --- 不是。前向使用 \texttt{sqrt\_approx\_fixpoint()}（Quadratic）或 C 标准库 \texttt{sqrt()}（Extended），逆向使用整数平方和移位。定点近似不可避免地引入舍入误差。
  \item \textbf{``$v_{\mathit{dco}}$ 总是 0''} --- 当 $\sigma = 0$ 且 $\alpha = 0$ 时确实为 0。但用户可以设置非零值来调整暗部映射。
  \item \textbf{``Extended 的 $T_1, T_2$ 和 Quadratic 的参数有关''} --- 无关。Extended 和 Quadratic 是完全独立的两种变换，参数不互通。
\end{enumerate}
\end{misunderstanding}

\section{与前后章衔接}

\begin{itemize}
  \item \textbf{输入}：本章的输入是原始图像像素值，存储在 \texttt{xs\_image\_t.comps\_array[]} 中
  \item \textbf{输出}：本章的输出（NLT 变换后的像素值）传递给第 \ref{ch:dwt} 章的前一步——第 05 章（MCT）做颜色空间变换
\end{itemize}

\section{本章小结}

NLT 是 JPEG XS 编码管线的第一步，负责将无符号像素值映射到以零为中心的有符号系数。三种模式从简到繁：None（纯线性）、Quadratic（全局二次压缩）、Extended（三段式灵活控制）。选择哪种模式取决于应用对压缩率和复杂度的要求。

下一章将讲解 MCT（多分量变换），它在 NLT 之后执行，负责利用分量间的相关性（色彩去相关）。
